The full Holzapfel-Ogden case was repeated with a thin wall ($R_e/R_i = 1.2$) and a very
thick wall ($R_e/R_i = 2.5$; the base case has 1.73), with the same inner radius, length
and material. The 3D mesh again has nr elements through the wall and about cubic elements,
so the thin wall needs more elements around and along the tube (26040 at nr = 4) and the
thick wall fewer (6480 at nr = 8). Both converge to the 0D block at the same rates as the
base case: the volume at orders 3.1--3.8, the radial displacement at 2.5--2.7, the
strains and stresses at 1.5--2.1. At the same nr the relative differences grow with the
wall thickness, since the elements are larger and the through-wall gradients steeper
(volume at nr = 4: $1.5\cdot10^{-7}$ thin, $1.5\cdot10^{-5}$ base, $1.8\cdot10^{-4}$ thick).
The nearly membrane-like thin wall agrees to $2\cdot10^{-5}$ in the volume already with one
element through the wall. The very thick wall reaches $1.1\cdot10^{-6}$ in the volume and
$4.7\cdot10^{-5}$ in $\rho$ at nr = 16 (48960 elements), where its 0D reference is still
accurate to 8\% and 1\% of these differences (Section~\ref{sec:zero-d}). Neither geometry
shows a model difference.
